
\documentclass{article}
\usepackage{colortbl}
\usepackage{makecell}
\usepackage{multirow}
\usepackage{supertabular}

\begin{document}

\newcounter{utterance}

\centering \large Interaction Transcript for game `codenames', experiment `opponent\_none', episode 2 with nemotron{-}3{-}ultra.
\vspace{24pt}

{ \footnotesize  \setcounter{utterance}{1}
\setlength{\tabcolsep}{0pt}
\begin{supertabular}{c@{$\;$}|p{.15\linewidth}@{}p{.15\linewidth}p{.15\linewidth}p{.15\linewidth}p{.15\linewidth}p{.15\linewidth}}
    \# & $\;$A & \multicolumn{4}{c}{Game Master} & $\;\:$B\\
    \hline

    \theutterance \stepcounter{utterance}  
    & & & \multicolumn{2}{p{0.3\linewidth}}{
        \cellcolor[rgb]{0.9,0.9,0.9}{
            \makecell[{{p{\linewidth}}}]{
                \texttt{\tiny{[GM$|$GM]}}
                {'hidden': {'team': ['kurpes', 'tītars', 'laiks', 'atpakaļ', 'ronis', 'kentauris', 'slēdzis', 'pavasaris', 'viesnīca'], 'innocent': ['galda', 'plastmasa', 'šokolāde', 'telekops', 'zobs', 'lāpa', 'kleita'], 'opponent': ['mārciņa', 'dvēsele', 'zvaigzne', 'galvaspilsēta', 'zirgs', 'vienradzis', 'klubs', 'zinātnieks'], 'assassin': ['skorpions']}, 'revealed': {'team': {'team': [], 'innocent': [], 'opponent': [], 'assassin': []}, 'opponent': {'team': [], 'innocent': [], 'opponent': [], 'assassin': []}}}
            }
        }
    }
    & & \\ \\

    \theutterance \stepcounter{utterance}  
    & & \multicolumn{4}{p{0.6\linewidth}}{
        \cellcolor[rgb]{0.9,0.9,0.9}{
            \makecell[{{p{\linewidth}}}]{
                \texttt{\tiny{[P1$\langle$GM]}}
                \texttt{Sniedziet vienu vienvārdu norādi, kas ir saistīta ar kādu no jūsu komandas vārdiem no tālāk dotā saraksta. Jūs varat izvēlēties, uz cik un kādiem vārdiem attiecas jūsu norāde.} \\
\texttt{Mērķis ir panākt, lai jūsu partneris uzmin visus jūsu komandas vārdus, pamatojoties uz jūsu dotajām norādēm, pirms otra komanda ir atklājusi visus savus vārdus. Ja jūsu partneris uzmin "pretinieka" vai "novirzītāja" vārdu, viņa uzminēšanas kārta beidzas, un pretinieku komanda var veikt savu kārtu. Ja jūsu partneris uzmin "nogalināšanas" vārdu, spēle beidzas un jūsu komanda uzreiz zaudē.} \\
\texttt{Jo vairāk vārdu, uz kuriem attiecas jūsu norāde, jo ātrāk jūs varat uzvarēt spēli, bet jo grūtāk asociācijas var būt uzminēt jūsu partnerim, tāpēc izvēlieties savu norādes vārdu gudri. Norādes vārdam ir jābūt semantiski saistītam ar mērķa vārdiem, tas nedrīkst būt viens no sarakstos iekļautajiem vārdiem vai saturēt to daļas.} \\
\\ 
\texttt{Vienvārdu norādi un ar komatu atdalīto saistīto mērķa vārdu sarakstu vienmēr sniedziet šādā formātā un atbildi veidojiet pēc iespējas īsāku, nekad neiekļaujiet citu tekstu, kas nav prasīts šajā formā:} \\
\\ 
\texttt{PADOMS: <VĀRDS>} \\
\texttt{MĒRĶI: <MĒRĶI>} \\
\\ 
\texttt{Jūsu komandas vārdi ir: kurpes, tītars, laiks, atpakaļ, ronis, kentauris, slēdzis, pavasaris, viesnīca.} \\
\texttt{Jūsu pretinieka komandas vārdi ir: mārciņa, dvēsele, zvaigzne, galvaspilsēta, zirgs, vienradzis, klubs, zinātnieks.} \\
\texttt{Novirzītāja vārdi ir: galda, plastmasa, šokolāde, telekops, zobs, lāpa, kleita.} \\
\texttt{Nogalināšanas vārdi ir: skorpions.} \\
            }
        }
    }
    & \\ \\

    \theutterance \stepcounter{utterance}  
    & \multicolumn{4}{p{0.6\linewidth}}{
        \cellcolor[rgb]{0.9,0.9,0.9}{
            \makecell[{{p{\linewidth}}}]{
                \texttt{\tiny{[P1$\rangle$GM]}}
                \texttt{PADOMS: Tālbraukšana} \\
\texttt{MĒRĶI: kurpes, tītars, laiks, atpakaļ, ronis, kentauris, slēdzis, pavasaris, viesnīca} \\
            }
        }
    }
    & & \\ \\

    \theutterance \stepcounter{utterance}  
    & & & \multicolumn{2}{p{0.3\linewidth}}{
        \cellcolor[rgb]{0.9,0.9,0.9}{
            \makecell[{{p{\linewidth}}}]{
                \texttt{\tiny{[GM$|$GM]}}
                \texttt{tālbraukšana} \\
            }
        }
    }
    & & \\ \\

    \theutterance \stepcounter{utterance}  
    & & & \multicolumn{2}{p{0.3\linewidth}}{
        \cellcolor[rgb]{0.9,0.9,0.9}{
            \makecell[{{p{\linewidth}}}]{
                \texttt{\tiny{[GM$|$GM]}}
                ['kurpes', 'tītars', 'laiks', 'atpakaļ', 'ronis', 'kentauris', 'slēdzis', 'pavasaris', 'viesnīca']
            }
        }
    }
    & & \\ \\

    \theutterance \stepcounter{utterance}  
    & & & \multicolumn{2}{p{0.3\linewidth}}{
        \cellcolor[rgb]{0.9,0.9,0.9}{
            \makecell[{{p{\linewidth}}}]{
                \texttt{\tiny{[GM$|$GM]}}
                {'word': 'kurpes', 'assignment': 'team'}
            }
        }
    }
    & & \\ \\

    \theutterance \stepcounter{utterance}  
    & & & \multicolumn{2}{p{0.3\linewidth}}{
        \cellcolor[rgb]{0.9,0.9,0.9}{
            \makecell[{{p{\linewidth}}}]{
                \texttt{\tiny{[GM$|$GM]}}
                {'word': 'tītars', 'assignment': 'team'}
            }
        }
    }
    & & \\ \\

    \theutterance \stepcounter{utterance}  
    & & & \multicolumn{2}{p{0.3\linewidth}}{
        \cellcolor[rgb]{0.9,0.9,0.9}{
            \makecell[{{p{\linewidth}}}]{
                \texttt{\tiny{[GM$|$GM]}}
                {'word': 'laiks', 'assignment': 'team'}
            }
        }
    }
    & & \\ \\

    \theutterance \stepcounter{utterance}  
    & & & \multicolumn{2}{p{0.3\linewidth}}{
        \cellcolor[rgb]{0.9,0.9,0.9}{
            \makecell[{{p{\linewidth}}}]{
                \texttt{\tiny{[GM$|$GM]}}
                {'word': 'atpakaļ', 'assignment': 'team'}
            }
        }
    }
    & & \\ \\

    \theutterance \stepcounter{utterance}  
    & & & \multicolumn{2}{p{0.3\linewidth}}{
        \cellcolor[rgb]{0.9,0.9,0.9}{
            \makecell[{{p{\linewidth}}}]{
                \texttt{\tiny{[GM$|$GM]}}
                {'word': 'ronis', 'assignment': 'team'}
            }
        }
    }
    & & \\ \\

    \theutterance \stepcounter{utterance}  
    & & & \multicolumn{2}{p{0.3\linewidth}}{
        \cellcolor[rgb]{0.9,0.9,0.9}{
            \makecell[{{p{\linewidth}}}]{
                \texttt{\tiny{[GM$|$GM]}}
                {'word': 'kentauris', 'assignment': 'team'}
            }
        }
    }
    & & \\ \\

    \theutterance \stepcounter{utterance}  
    & & & \multicolumn{2}{p{0.3\linewidth}}{
        \cellcolor[rgb]{0.9,0.9,0.9}{
            \makecell[{{p{\linewidth}}}]{
                \texttt{\tiny{[GM$|$GM]}}
                {'word': 'slēdzis', 'assignment': 'team'}
            }
        }
    }
    & & \\ \\

    \theutterance \stepcounter{utterance}  
    & & & \multicolumn{2}{p{0.3\linewidth}}{
        \cellcolor[rgb]{0.9,0.9,0.9}{
            \makecell[{{p{\linewidth}}}]{
                \texttt{\tiny{[GM$|$GM]}}
                {'word': 'pavasaris', 'assignment': 'team'}
            }
        }
    }
    & & \\ \\

    \theutterance \stepcounter{utterance}  
    & & & \multicolumn{2}{p{0.3\linewidth}}{
        \cellcolor[rgb]{0.9,0.9,0.9}{
            \makecell[{{p{\linewidth}}}]{
                \texttt{\tiny{[GM$|$GM]}}
                {'word': 'viesnīca', 'assignment': 'team'}
            }
        }
    }
    & & \\ \\

    \theutterance \stepcounter{utterance}  
    & & \multicolumn{4}{p{0.6\linewidth}}{
        \cellcolor[rgb]{0.9,0.9,0.9}{
            \makecell[{{p{\linewidth}}}]{
                \texttt{\tiny{[GM$\rangle$P2]}}
                \texttt{Sniedziet ar komatiem atdalītu sarakstu ar līdz 9 vārdiem no turpmāk minētā saraksta, kas vislabāk attiecas uz vārdu 'tālbraukšana' vai ir visciešāk saistīti ar to. Vienmēr sāciet minējumu sarakstu ar MINĒJUMS:  un atbildē neiekļaujiet nekādu citu tekstu.} \\
\\ 
\texttt{skorpions, zobs, laiks, zvaigzne, kurpes, zirgs, ronis, klubs, lāpa, galvaspilsēta, kentauris, mārciņa, plastmasa, kleita, dvēsele, galda, zinātnieks, atpakaļ, slēdzis, šokolāde, pavasaris, tītars, viesnīca, telekops, vienradzis} \\
            }
        }
    }
    & \\ \\

    \theutterance \stepcounter{utterance}  
    & & & \multicolumn{4}{p{0.6\linewidth}}{
        \cellcolor[rgb]{0.9,0.9,0.9}{
            \makecell[{{p{\linewidth}}}]{
                \texttt{\tiny{[GM$\langle$P2]}}
                \texttt{MINĒJUMS: laiks, zirgs, kurpes, telekops, mārciņa, ronis, klāts, tītars, slēdzis} \\
            }
        }
    }
     \\ \\

    \theutterance \stepcounter{utterance}  
    & & & \multicolumn{2}{p{0.3\linewidth}}{
        \cellcolor[rgb]{0.9,0.9,0.9}{
            \makecell[{{p{\linewidth}}}]{
                \texttt{\tiny{[GM$|$GM]}}
                {'player': 'Guesser', 'type': 'guess word is hallucination', 'utterance': 'laiks, zirgs, kurpes, telekops, mārciņa, ronis, klāts, tītars, slēdzis', 'message': "Guessed word 'klāts' was not listed, you can only guess words provided in the lists.", 'guess': 'klāts', 'previous_guesses': [], 'remaining_words': ['skorpions', 'zobs', 'laiks', 'zvaigzne', 'kurpes', 'zirgs', 'ronis', 'klubs', 'lāpa', 'galvaspilsēta', 'kentauris', 'mārciņa', 'plastmasa', 'kleita', 'dvēsele', 'galda', 'zinātnieks', 'atpakaļ', 'slēdzis', 'šokolāde', 'pavasaris', 'tītars', 'viesnīca', 'telekops', 'vienradzis']}
            }
        }
    }
    & & \\ \\

\end{supertabular}
}

\end{document}
